
\subsection{余导子与 Hochschild 上同调}
\label{subsec:4-4-5}

现在继续证明定理 4.2.34 在滤过 $A_\infty$ 代数情形下的结论。证明
将对能量滤过作归纳。\liftIndex{I164}

先引入一些记号。令 $\overline C_i$（$i=1,2$）为自由分次 $R$-模。
自由分次 $R$-模
$B\overline C_i[1]=\bigoplus_{k=0}^{\infty}B_k\overline C_i[1]$
带有由 (3.2.14) 给出的余乘法 $\Delta$。令
$\widehat{\overline{\mathfrak f}}\colon
B\overline C_1[1]\to B\overline C_2[1]$ 为次数 $0$ 的余代数同态，
并置
\[
 \varphi=(\varphi_0,\varphi_1,\varphi_2,\ldots)
 \in\prod_{k=0}^{\infty}
 \operatorname{Hom}(B_k\overline C_1[1],\overline C_2[1]),
\]
其中 $\deg\varphi_k=\deg\varphi$ 与 $k$ 无关。以下按照约定 3
中的 (1.35)，定义
$\widehat\varphi\colon B\overline C_1[1]\to B\overline C_2[1]$：
\begin{equation}
 \widehat\varphi(\mathbf x)
 =\sum_a(-1)^{\deg\varphi\,\deg'\mathbf x_a^{(3;1)}}
 \widehat{\overline{\mathfrak f}}(\mathbf x_a^{(3;1)})
 \widehat\otimes\varphi(\mathbf x_a^{(3;2)})
 \widehat\otimes
 \widehat{\overline{\mathfrak f}}(\mathbf x_a^{(3;3)}).
\tag{4.4.39}\label{eq:4-4-39}
\end{equation}
这里 $\varphi\colon B\overline C_1[1]\to\overline C_2[1]$ 是在
$B_k\overline C_1[1]$ 上限制为 $\varphi_k$ 的同态。容易验证，
$\widehat\varphi$ 是一个余导子\liftIndex{I163}，即满足
\begin{equation}
 \bigl(\widehat\varphi\widehat\otimes
 \widehat{\overline{\mathfrak f}}+
 \widehat{\overline{\mathfrak f}}\widehat\otimes
 \widehat\varphi\bigr)\circ\Delta
 =\Delta\circ\widehat\varphi.
\tag{4.4.40}\label{eq:4-4-40}
\end{equation}
此处 $\widehat\otimes$ 是分次张量积，例如
\[
 (\widehat{\overline{\mathfrak f}}\widehat\otimes
 \widehat\varphi)(\mathbf x\otimes\mathbf y)
 =(-1)^{\deg\varphi\,\deg'\mathbf x}
 \widehat{\overline{\mathfrak f}}(\mathbf x)
 \otimes\widehat\varphi(\mathbf y).
\]
还要注意，$\widehat\varphi$ 的像落在
$B_{1\cdots\infty}\overline C_2[1]$ 中；换言之，它的
$B_0\overline C_2[1]=R$ 分量为零。置
\[
 B_{1\cdots\infty}\overline C_i[1]
 =\bigoplus_{k=1}^{\infty}B_k\overline C_i[1].
\]

\begin{definition}{4.4.41}\label{def:4-4-41}
记
$\operatorname{Der}(B\overline C_1[1],B\overline C_2[1];
\overline{\mathfrak f})$
为所有像落在 $B_{1\cdots\infty}\overline C_2[1]$ 中的分次余导子
$\widehat\varphi$ 所组成的集合。
\end{definition}

映射 $\varphi\mapsto\widehat\varphi$ 给出同态
\begin{equation}
 \prod_{k=0}^{\infty}
 \operatorname{Hom}(B_k\overline C_1[1],\overline C_2[1])
 \longrightarrow
 \operatorname{Der}(B\overline C_1[1],B\overline C_2[1];
 \overline{\mathfrak f}).
\tag{4.4.42}\label{eq:4-4-42}
\end{equation}
下面证明：

\begin{lemma}{4.4.43}\label{lem:4-4-43}
式 \eqref{eq:4-4-42} 是同构。\translatornote{原书在本引理证明中，
把分量 $\varphi_k$ 的普通 $\operatorname{Hom}$ 空间写作
“$\operatorname{Hom}(\ldots;\overline{\mathfrak f})$”；式
(4.4.42)、(4.4.44) 的直积均与 $\overline{\mathfrak f}$ 无关，
依赖它的是余导子空间，故译文不把此参数附在普通同态空间上。}
\end{lemma}

\begin{proof}
令 $\widehat\varphi\in
\operatorname{Der}(B\overline C_1[1],B\overline C_2[1];
\overline{\mathfrak f})$，取其分量
$\varphi_k\in\operatorname{Hom}
(B_k\overline C_1[1],B_1\overline C_2[1])$。这样得到同态
\begin{equation}
 \operatorname{Der}(B\overline C_1[1],B\overline C_2[1];
 \overline{\mathfrak f})
 \longrightarrow
 \prod_{k=0}^{\infty}
 \operatorname{Hom}(B_k\overline C_1[1],\overline C_2[1]).
\tag{4.4.44}\label{eq:4-4-44}
\end{equation}
容易看出，先作 \eqref{eq:4-4-42} 再作 \eqref{eq:4-4-44}，所得
复合是该直积上的恒等映射。因此只需证明 \eqref{eq:4-4-44} 是单射。

设 $\widehat\varphi$ 落在 \eqref{eq:4-4-44} 的核中。令
$\pi_{1\cdots n}\colon B\overline C_2[1]\to
B_{1\cdots n}\overline C_2[1]$ 为投影。将对 $n$ 归纳证明
$\pi_{1\cdots n}\circ\widehat\varphi=0$。注意，由定义 4.4.41，
$\pi_0\circ\widehat\varphi=0$。

$n=1$ 的情形直接来自 $\widehat\varphi$ 落在
\eqref{eq:4-4-44} 的核中。假设
$\pi_{1\cdots n}\circ\widehat\varphi=0$，现证明
$\pi_{1\cdots n+1}\circ\widehat\varphi=0$。令
$\pi_k\colon B\overline C_i[1]\to B_k\overline C_i[1]$ 为投影。
归纳假设给出
\begin{equation}
 (\pi_{1\cdots\infty}\otimes\pi_{1\cdots\infty})
 \circ\Delta\circ\pi_{1\cdots n+1}\circ\widehat\varphi
 =\sum_{k=1}^{n}(\pi_k\otimes\pi_{n+1-k})
 \circ\Delta\circ\widehat\varphi.
\tag{4.4.45}\label{eq:4-4-45}
\end{equation}
另一方面，
\[
 \begin{aligned}
 &\sum_{k=1}^{n}(\pi_k\otimes\pi_{n+1-k})
 \circ\Delta\circ\widehat\varphi\\
 &\qquad=
 \sum_{k=1}^{n}(\pi_k\otimes\pi_{n+1-k})\circ
 \bigl((\widehat\varphi\widehat\otimes
 \widehat{\overline{\mathfrak f}})
 +(\widehat{\overline{\mathfrak f}}\widehat\otimes
 \widehat\varphi)\bigr)\circ\Delta.
 \end{aligned}
\]
由归纳假设，右端为零。对 $B_{n+1}\overline C_2[1]$ 而言，
$(\pi_1\otimes\pi_n)\circ\Delta$ 是把一个长度 $n+1$ 的张量
在第一个因子之后切开的自然同构。故上式的 $k=1$ 分量为零便推出
$\pi_{n+1}\circ\widehat\varphi=0$，完成归纳，引理得证。
\translatornote{原书在此改取余积的 $B_0\otimes B_{n+1}$ 与
$B_{n+1}\otimes B_0$ 分量，并据 $\pi_0\widehat\varphi=0$ 宣称它们为零；
但式 (3.2.14) 的余积含空张量项，此结论不能推出。译文改用已经算得为零
的 $B_1\otimes B_n$ 分量完成归纳。}
\end{proof}

现在再假设 $\overline C_i$ 具有 $A_\infty$ 代数结构，令
$\widehat{\overline d}$ 为 $B\overline C_i[1]$ 上的上边界算子。
还假设 $\widehat{\overline{\mathfrak f}}$ 由一个 $A_\infty$ 同态
$\overline{\mathfrak f}$ 诱导。于是
$\widehat{\overline{\mathfrak f}}$ 是一个
$\widehat{\overline d}$-上链映射。

\begin{statement}{引理--定义}{4.4.46}\label{lem-def:4-4-46}
若
$\widehat\varphi\in
\operatorname{Der}(B\overline C_1[1],B\overline C_2[1];
\overline{\mathfrak f})$，则
\[
 \widehat d\circ\widehat\varphi
 -(-1)^{\deg\widehat\varphi}
 \widehat\varphi\circ\widehat d
 \in\operatorname{Der}(B\overline C_1[1],B\overline C_2[1];
 \overline{\mathfrak f}).
\]
据此定义
\[
 \delta\colon
 \operatorname{Der}(B\overline C_1[1],B\overline C_2[1];
 \overline{\mathfrak f})
 \longrightarrow
 \operatorname{Der}(B\overline C_1[1],B\overline C_2[1];
 \overline{\mathfrak f}),
\]
其作用为
$\delta(\widehat\varphi)=
\widehat d\circ\widehat\varphi
-(-1)^{\deg\widehat\varphi}\widehat\varphi\circ\widehat d$。
\end{statement}

证明是标准计算，此处从略。容易看出 $\delta^2=0$。

\begin{definition}{4.4.47}\label{def:4-4-47}
定义
\[
 HH(\overline C_1,\overline C_2;\overline{\mathfrak f})
 :=\ker\delta/\operatorname{Im}\delta,
\]
称为 $A_\infty$ 同态 $\overline{\mathfrak f}$ 的
\emph{Hochschild 上同调}\liftIndex{I165}。
\end{definition}

注意，\eqref{eq:4-4-42} 的左端与 $\overline{\mathfrak f}$ 无关，
右端却依赖于它。因此，Hochschild 上同调一般也依赖于
$\overline{\mathfrak f}$。

令
$\overline{\mathfrak g}\colon\overline C'_1\to\overline C_1$、
$\overline{\mathfrak g}'\colon\overline C_2\to\overline C'_2$
为 $A_\infty$ 同态。它们以自然方式诱导上链映射
\[
 (\overline{\mathfrak g},\overline{\mathfrak g}')_*
 \colon
 \operatorname{Der}(B\overline C_1[1],B\overline C_2[1];
 \overline{\mathfrak f})
 \longrightarrow
 \operatorname{Der}(B\overline C'_1[1],B\overline C'_2[1];
 \overline{\mathfrak g}'\circ\overline{\mathfrak f}
 \circ\overline{\mathfrak g}).
\]
它在 Hochschild 上同调上诱导的映射仍记为
$(\overline{\mathfrak g},\overline{\mathfrak g}')_*$。

\begin{proposition}{4.4.48}\label{prop:4-4-48}
若 $\overline{\mathfrak g}$ 与 $\overline{\mathfrak g}'$ 都是同伦等价，
则
\[
 (\overline{\mathfrak g},\overline{\mathfrak g}')_*
 \colon HH(\overline C_1,\overline C_2;\overline{\mathfrak f})
 \longrightarrow
 HH(\overline C'_1,\overline C'_2;
 \overline{\mathfrak g}'\circ\overline{\mathfrak f}
 \circ\overline{\mathfrak g})
\]
是同构。\liftIndex{I166}
\end{proposition}

\begin{proof}
令 $\overline{\mathfrak C}_1$、$\overline{\mathfrak C}_2$ 分别为
$[0,1]\times\overline C_1$、$[0,1]\times\overline C_2$ 的模型。
$\overline{\mathrm{Incl}}$ 诱导同态
\[
 (\overline{\mathrm{Incl}},\operatorname{id})_*
 \colon HH(\overline{\mathfrak C}_1,\overline C_2)
 \longrightarrow HH(\overline C_1,\overline C_2)
\]
以及
\[
 (\operatorname{id},\overline{\mathrm{Incl}})_*
 \colon HH(\overline C_1,\overline C_2)
 \longrightarrow HH(\overline C_1,\overline{\mathfrak C}_2).
\]

\begin{lemma}{4.4.49}\label{lem:4-4-49}
$(\overline{\mathrm{Incl}},\operatorname{id})_*$ 与
$(\operatorname{id},\overline{\mathrm{Incl}})_*$ 都是同构。
\end{lemma}

\begin{proof}
先证 $(\overline{\mathrm{Incl}},\operatorname{id})_*$ 的情形。由于
\[
 (\overline{\mathrm{Incl}},\operatorname{id})_*
 \circ(\overline{\mathrm{Eval}}_{s=0},\operatorname{id})_*
 =\operatorname{id},
\]
它是满射。现证其单射性。令
\[
 \varphi=(\varphi_0,\varphi_1,\ldots)
 \in\prod_k\operatorname{Hom}
 (B_k\overline{\mathfrak C}_1[1],\overline C_2[1])
\]
满足 $\delta\varphi=0$，且
$\varphi\circ\widehat{\overline{\mathrm{Incl}}}=\delta\psi$。
则
\[
 \delta\bigl(\varphi-
 \delta(\psi\circ\widehat{\overline{\mathrm{Eval}}}_{s=0})\bigr)=0,
\]
并且
\[
 \bigl(\varphi-
 \delta(\psi\circ\widehat{\overline{\mathrm{Eval}}}_{s=0})\bigr)
 \circ\widehat{\overline{\mathrm{Incl}}}=0.
\]
因此只须考虑满足
\begin{equation}
 \delta\varphi=0,
 \qquad
 \varphi\circ\widehat{\overline{\mathrm{Incl}}}=0
\tag{4.4.50}\label{eq:4-4-50}
\end{equation}
的 $\varphi$，并证明下面的子引理。

\begin{statement}{子引理}{4.4.51}\label{sublem:4-4-51}
若 $\varphi=(\varphi_0,\varphi_1,\ldots)$ 满足
\eqref{eq:4-4-50}，且
\[
 \varphi\in\prod_{\ell=k}^{\infty}
 \operatorname{Hom}(B_\ell\overline{\mathfrak C}_1[1],
 \overline C_2[1]),
\]
则存在
\[
 \psi\in\prod_{\ell=k}^{\infty}
 \operatorname{Hom}(B_\ell\overline{\mathfrak C}_1[1],
 \overline C_2[1]),
\]
使得
\[
 \varphi-\delta\psi\in\prod_{\ell=k+1}^{\infty}
 \operatorname{Hom}(B_\ell\overline{\mathfrak C}_1[1],
 \overline C_2[1]),
\]
而且 $\varphi-\delta\psi$ 仍满足 \eqref{eq:4-4-50}。
\end{statement}

\begin{proof}
注意
\[
 \widehat{\overline{\mathrm{Incl}}}\colon
 \operatorname{Hom}(B_k\overline{\mathfrak C}_1[1],
 \overline C_2[1])
 \longrightarrow
 \operatorname{Hom}(B_k\overline C_1[1],
 \overline C_2[1])
\]
满足
$\delta_1\circ\widehat{\overline{\mathrm{Incl}}}
=\widehat{\overline{\mathrm{Incl}}}\circ\delta_1$，
且在 $\delta_1$-上同调上诱导同构。这里
$\delta_1(A)=\overline{\mathfrak m}_1\circ A
-(-1)^{\deg A}A\circ\widehat{\overline{\mathfrak m}}_1$。

由于 $\delta\varphi=0$，且 $\varphi$ 没有低于 $k$ 的分量，
可得 $\delta_1(\varphi_k)=0$。又由假设
$\varphi_k\circ\widehat{\overline{\mathrm{Incl}}}=0$，故存在
$\psi'_k\in\operatorname{Hom}
(B_k\overline{\mathfrak C}_1[1],\overline C_2[1])$，使
$\varphi_k=\delta_1(\psi'_k)$。置
\[
 \psi'=(0,\ldots,0,\psi'_k,0,\ldots)
 \in\prod_{\ell=k}^{\infty}
 \operatorname{Hom}(B_\ell\overline{\mathfrak C}_1[1],
 \overline C_2[1]),
\]
并令
\[
 \psi=\psi'-\psi'\circ
 \widehat{\overline{\mathrm{Incl}}}
 \circ\widehat{\overline{\mathrm{Eval}}}_{s=0}.
\]
直接计算即可验证 $\psi$ 具有所需性质。
\translatornote{原书在此写“$\widehat\varphi=0$”，这与待证的非零闭上链
$\varphi$ 矛盾；由子引理的假设 (4.4.50) 及随后的推论
$\delta_1(\varphi_k)=0$，此处应为 $\delta\varphi=0$。此外原书三处用
直和记载可有无限多高次分量的余导子；译文依式 (4.4.42) 的直积
及逐阶归纳改为直积。}
\end{proof}

由子引理 4.4.51，
$(\overline{\mathrm{Incl}},\operatorname{id})_*$ 的结论立即成立。
再证 $(\operatorname{id},\overline{\mathrm{Incl}})_*$。由于
\[
 (\operatorname{id},\overline{\mathrm{Eval}}_{s=0})_*
 \circ(\operatorname{id},\overline{\mathrm{Incl}})_*
 =\operatorname{id},
\]
它是单射，只需证明满射。令
\[
 \psi\in\prod_k\operatorname{Hom}
 (B_k\overline C_1[1],\overline{\mathfrak C}_2[1])
\]
为一个 $\delta$-闭上链，即 $\delta\psi=0$。更具体地说，
$\overline{\mathfrak m}_1\circ\psi
-(-1)^{\deg\psi}\psi\circ\widehat d=0$。注意
$\overline{\mathrm{Eval}}_{s=0}\circ\psi$
是
$\prod_k\operatorname{Hom}
(B_k\overline C_1[1],\overline C_2[1])$
中的 $\delta$-闭上链。由于
$\overline{\mathrm{Incl}}\circ
\overline{\mathrm{Eval}}_{s=0}$ 与恒等映射上链同伦，即
\[
 \operatorname{id}-\overline{\mathrm{Incl}}\circ
 \overline{\mathrm{Eval}}_{s=0}
 =\overline{\mathfrak m}_1\circ h
 +h\circ\overline{\mathfrak m}_1,
\]
其中 $h$ 是上链同伦，因此
\[
 \begin{aligned}
 &\psi-\overline{\mathrm{Incl}}\circ
 \overline{\mathrm{Eval}}_{s=0}\circ\psi
 =(\overline{\mathfrak m}_1\circ h
 +h\circ\overline{\mathfrak m}_1)\psi\\
 &\qquad=\overline{\mathfrak m}_1(h\circ\psi)
 +(-1)^{\deg\psi}(h\circ\psi)\circ\widehat d\\
 &\qquad=\overline{\mathfrak m}_1(h\circ\psi)
 -(-1)^{\deg(h\circ\psi)}
 (h\circ\psi)\circ\widehat d.
 \end{aligned}
\]
也就是说，
\[
 \psi=\overline{\mathrm{Incl}}
 (\overline{\mathrm{Eval}}_{s=0}\circ\psi)
 +\delta(h\circ\psi).
\]
故 $(\operatorname{id},\overline{\mathrm{Incl}})_*$ 为满射。引理
4.4.49 得证。
\end{proof}

利用引理 4.4.49，容易证明
$(\overline{\mathfrak g},\overline{\mathfrak g}')_*$
只依赖于 $\overline{\mathfrak g}$ 与
$\overline{\mathfrak g}'$ 的同伦类。这就证明了命题 4.4.48。
\end{proof}

\subsection{能量滤过归纳}
\label{subsec:4-4-6}

本小节及下一小节完成定理 4.2.34 的证明，仍采用该定理中的记号。\liftIndex{I167}
取满足条件 3.1.6 的 $G\subset\mathbb R_{\geq0}\times2\mathbb Z$，
并假定所考虑的滤过 $A_\infty$ 代数及其同态均为 $G$-有隙的。
记 $G$ 到 $\mathbb R_{\geq0}$ 的投影之像为
$\operatorname{Spec}(G)=\{\lambda_0,\lambda_1,\ldots\}$。
由条件 3.1.6，$\operatorname{Spec}(G)$ 是离散的，因此可依次排列，
使 $\lambda_i<\lambda_{i+1}$ 且
$\lim_{i\to\infty}\lambda_i=\infty$。我们将构造
\[
 \widetilde{\mathfrak F}_{k,i}\colon
 B_k\overline{\mathfrak C}_1[1]
 \longrightarrow\overline{\mathfrak C}_2[1],
 \qquad k=0,1,2,\ldots,\quad i=0,1,2,\ldots,
\]
并定义
\[
 \widetilde{\mathfrak F}^{(n)}_k
 =\sum_{i=0}^{n}\widetilde{\mathfrak F}_{k,i}T^{\lambda_i}
 \colon B_k\overline{\mathfrak C}_1[1]
 \longrightarrow\overline{\mathfrak C}_2[1],
 \qquad k=0,1,2,\ldots.
\]
这里假设 $\widetilde{\mathfrak F}_{0,0}=0$。把
$\sum_k\widetilde{\mathfrak F}^{(n)}_k$ 延拓为余代数同态
$\widehat{\widetilde{\mathfrak F}}^{(n)}\colon
\widehat B\mathfrak C_1[1]\to\widehat B\mathfrak C_2[1]$。
下面对 $n$ 归纳证明：

\begin{proposition}{4.4.52}\label{prop:4-4-52}
对任意 $n=0,1,2,\ldots$，存在 $i\leq n$ 的
$\widetilde{\mathfrak F}_{k,i}$，使得
\begin{align}
 \widehat d\circ\widehat{\widetilde{\mathfrak F}}^{(n)}
 &\equiv\widehat{\widetilde{\mathfrak F}}^{(n)}
 \circ\widehat d\pmod{T^{\lambda_{n+1}}},
 \tag{4.4.53.1}\label{eq:4-4-53-1}\\
 \widehat{\mathrm{Incl}}\circ\widehat{\mathfrak f}
 &\equiv\widehat{\widetilde{\mathfrak F}}^{(n)}
 \circ\widehat{\mathrm{Incl}}\pmod{T^{\lambda_{n+1}}},
 \tag{4.4.53.2}\label{eq:4-4-53-2}\\
 \widehat{\mathrm{Eval}}_{s=s_0}
 \circ\widehat{\widetilde{\mathfrak F}}^{(n)}
 &\equiv\widehat{\mathfrak f}
 \circ\widehat{\mathrm{Eval}}_{s=s_0}
 \pmod{T^{\lambda_{n+1}}},\quad s_0=0,1.
 \tag{4.4.53.3}\label{eq:4-4-53-3}
\end{align}
此外，
\begin{equation}
 \widehat{\widetilde{\mathfrak F}}^{(n)}
 -\widehat{\widetilde{\mathfrak F}}^{(n-1)}
 \equiv0\pmod{T^{\lambda_n}}.
\tag{4.4.53.4}\label{eq:4-4-53-4}
\end{equation}
\end{proposition}

\begin{proof}
把定理 4.2.34 的无滤过版本用于
$\overline{\mathfrak f}$，得到 $\overline{\mathfrak F}_k$。
置 $\widetilde{\mathfrak F}_{k,0}=\overline{\mathfrak F}_k$，即得
$n=0$ 的情形。

假设命题对 $n$ 成立，现证明对 $n+1$ 成立。由归纳假设
(4.4.53)，存在
\[
 \begin{gathered}
 \operatorname{Err1}_n\in
 \operatorname{Der}(B\overline{\mathfrak C}_1[1],
 B\overline{\mathfrak C}_2[1]),\\
 \operatorname{Err2}_n\in
 \operatorname{Der}(B\overline C_1[1],
 B\overline{\mathfrak C}_2[1]),\\
 \operatorname{Err3}_n\in
 \operatorname{Der}(B\overline{\mathfrak C}_1[1],
 B\overline C_2[1]),
 \end{gathered}
\]
使得
\begin{align}
 \widehat d\circ\widehat{\widetilde{\mathfrak F}}^{(n)}
 -\widehat{\widetilde{\mathfrak F}}^{(n)}\circ\widehat d
 -T^{\lambda_{n+1}}\operatorname{Err1}_n
 &\equiv0\pmod{T^{\lambda_{n+2}}},
 \tag{4.4.54.1}\label{eq:4-4-54-1}\\
 \widehat{\mathrm{Incl}}\circ\widehat{\mathfrak f}
 -\widehat{\widetilde{\mathfrak F}}^{(n)}
 \circ\widehat{\mathrm{Incl}}
 -T^{\lambda_{n+1}}\operatorname{Err2}_n
 &\equiv0\pmod{T^{\lambda_{n+2}}},
 \tag{4.4.54.2}\label{eq:4-4-54-2}\\
 \widehat{\mathrm{Eval}}_{s=s_0}
 \circ\widehat{\widetilde{\mathfrak F}}^{(n)}
 -\widehat{\mathfrak f}\circ
 \widehat{\mathrm{Eval}}_{s=s_0}
 -T^{\lambda_{n+1}}\operatorname{Err3}_n
 &\equiv0\pmod{T^{\lambda_{n+2}}}.
 \tag{4.4.54.3}\label{eq:4-4-54-3}
\end{align}
注意
\begin{equation}
 \widehat d\equiv\widehat{\overline d}
 \pmod{T^{\lambda_1}}.
\tag{4.4.55}\label{eq:4-4-55}
\end{equation}
由于 $\overline{\mathrm{Incl}}$ 是分裂单射，可以取
$\widetilde{\operatorname{Err2}}_n\in
\operatorname{Der}(B\overline{\mathfrak C}_1[1],
B\overline{\mathfrak C}_2[1])$，使
$\widetilde{\operatorname{Err2}}_n
\circ\widehat{\overline{\mathrm{Incl}}}
=\operatorname{Err2}_n$。将
$\widehat{\widetilde{\mathfrak F}}^{(n)}$ 替换为
$\widehat{\widetilde{\mathfrak F}}^{(n)}
+T^{\lambda_{n+1}}\widetilde{\operatorname{Err2}}_n$。\translatornote{原书在此把修正写成减去
$T^{\lambda_{n+1}}\widetilde{\operatorname{Err2}}_n$；但按
(4.4.54.2) 的误差定义及其与 $\widehat{\overline{\mathrm{Incl}}}$
复合后的等式，只有加上这一项才能得到 (4.4.54.2+)，译文据此改正。}
必要时调整 $\operatorname{Err1}_n$ 与
$\operatorname{Err3}_n$ 的选取，便可保持
\eqref{eq:4-4-54-1}、\eqref{eq:4-4-54-3}，并令
\begin{equation}
 \widehat{\mathrm{Incl}}\circ\widehat{\mathfrak f}
 -\widehat{\widetilde{\mathfrak F}}^{(n)}
 \circ\widehat{\mathrm{Incl}}
 \equiv0\pmod{T^{\lambda_{n+2}}}.
\tag{4.4.54.2+}\label{eq:4-4-54-2-plus}
\end{equation}
由 \eqref{eq:4-4-54-2-plus} 与 \eqref{eq:4-4-54-3} 得到
\begin{equation}
 \operatorname{Err3}_n
 \circ\widehat{\overline{\mathrm{Incl}}}=0
 \in\operatorname{Der}(B\overline C_1[1],B\overline C_2[1]).
\tag{4.4.56}\label{eq:4-4-56}
\end{equation}
利用 $(\overline{\mathrm{Eval}}_{s=0})_1$、
$(\overline{\mathrm{Eval}}_{s=1})_1$ 联合满射以及
$\overline{\mathrm{Incl}}_1$ 分裂单射的性质，并应用
\eqref{eq:4-4-56} 与引理 4.4.43，可找到余导子
$\widetilde{\operatorname{Err3}}_n\in
\operatorname{Der}(B\overline{\mathfrak C}_1[1],
B\overline{\mathfrak C}_2[1])$，使
\begin{align}
 \bigl(\widehat{\overline{\mathrm{Eval}}}_{s=0}
 \circ\widetilde{\operatorname{Err3}}_n,
 \widehat{\overline{\mathrm{Eval}}}_{s=1}
 \circ\widetilde{\operatorname{Err3}}_n\bigr)
 &=\operatorname{Err3}_n,
 \tag{4.4.57.1}\label{eq:4-4-57-1}\\
 \widetilde{\operatorname{Err3}}_n
 \circ\widehat{\overline{\mathrm{Incl}}}&=0.
 \tag{4.4.57.2}\label{eq:4-4-57-2}
\end{align}
将 $\widehat{\widetilde{\mathfrak F}}^{(n)}$ 替换为
$\widehat{\widetilde{\mathfrak F}}^{(n)}
-T^{\lambda_{n+1}}\widetilde{\operatorname{Err3}}_n$。
此时 \eqref{eq:4-4-54-1}、\eqref{eq:4-4-54-2-plus} 仍成立，且
\begin{equation}
 \widehat{\mathrm{Eval}}_{s=s_0}
 \circ\widehat{\widetilde{\mathfrak F}}^{(n)}
 -\widehat{\mathfrak f}\circ
 \widehat{\mathrm{Eval}}_{s=s_0}
 \equiv0\pmod{T^{\lambda_{n+2}}}.
\tag{4.4.54.3+}\label{eq:4-4-54-3-plus}
\end{equation}

现在有下面的引理。

\begin{lemma}{4.4.58}\label{lem:4-4-58}
$\delta(\operatorname{Err1}_n)=0$，这里的 $\delta$ 如定义
4.4.46 所述。
\end{lemma}

\begin{proof}
由定义，
\[
 \widehat d\circ\operatorname{Err1}_n
 +\operatorname{Err1}_n\circ\widehat d=0.
\]
结合 \eqref{eq:4-4-55} 即得结论。
\end{proof}

由 \eqref{eq:4-4-54-2-plus}、\eqref{eq:4-4-54-3-plus}，还有
\begin{align}
 \operatorname{Err1}_n
 \circ\widehat{\overline{\mathrm{Incl}}}_1&=0,
 \tag{4.4.59.1}\label{eq:4-4-59-1}\\
 \widehat{\overline{\mathrm{Eval}}}_{s=s_0}
 \circ\operatorname{Err1}_n&=0.
 \tag{4.4.59.2}\label{eq:4-4-59-2}
\end{align}

置
\[
 \overline D_3
 =\ker\overline{\mathrm{Eval}}_{0*}
 \cap\ker\overline{\mathrm{Eval}}_{1*}
 \subseteq\overline{\mathfrak C}_2[1].
\]
将引理 4.4.3 与引理 4.4.49 用于
$A=\operatorname{Err1}_n$，可得到
\[
 \widetilde{\mathfrak F}_{k,n+1}
 \in\operatorname{Hom}
 (B_k\overline{\mathfrak C}_1[1],\overline D_3)
 \subseteq
 \operatorname{Der}(B\overline{\mathfrak C}_1[1],
 B\overline{\mathfrak C}_2[1]),
 \qquad k=0,1,2,\ldots,
\]
使得
\begin{align}
 \widetilde{\mathfrak F}_{*,n+1}
 \circ\widehat{\overline{\mathrm{Incl}}}_1&=0,
 \tag{4.4.60.1}\label{eq:4-4-60-1}\\
 \delta(\widetilde{\mathfrak F}_{*,n+1})
 +\operatorname{Err1}_n&=0.
 \tag{4.4.60.2}\label{eq:4-4-60-2}
\end{align}
容易看出 $\widetilde{\mathfrak F}_{k,n+1}$ 具有所需性质。
命题 4.4.52 得证。
\end{proof}

由命题 4.4.52 推出定理 4.2.34 时，只须在
$\mathfrak f$ 为严格同态和／或含幺同态的情形下作补充论证。
若 $\mathfrak f$ 严格，可按如下方式归纳选取
$\widetilde{\mathfrak F}_{k,i}$，使
$\widetilde{\mathfrak F}_{0,i}=0$。假设 $i\leq n$ 时已有
$\widetilde{\mathfrak F}_{0,i}=0$。因为 $\mathfrak f$ 严格，
$\operatorname{Err1}_n$、$\operatorname{Err2}_n$、
$\operatorname{Err3}_n$ 在相应的
$B_0\overline C_1$ 或 $B_0\overline{\mathfrak C}_1$ 上为零。
因此，在修正 $\widehat{\widetilde{\mathfrak F}}^{(n)}$ 以满足
\eqref{eq:4-4-54-2-plus} 和 \eqref{eq:4-4-54-3-plus} 时，仍可保持
$\widetilde{\mathfrak F}_{0,i}=0$。于是
$\operatorname{Err1}_n$ 在 $B_0\overline{\mathfrak C}_1$ 上为零；
构造时即可取 $\widetilde{\mathfrak F}_{0,n+1}=0$。
含幺情形在下一小节处理。

\subsection{含幺情形 II：滤过版本}
\label{subsec:4-4-7}

本小节证明命题 4.4.52 的含幺滤过版本。令 $(C,\mathfrak m)$ 为
含幺滤过 $A_\infty$ 代数。

\begin{definition}{4.4.61}\label{def:4-4-61}
如下定义\emph{部分约化 bar 复形}
$\widehat B^{\mathrm{red}}C[1]$\liftIndex{I168}。置
$C_{\mathrm{red}}=C/\Lambda_{0,nov}[e]$，并定义
\begin{equation}
 B_k^{\mathrm{red}}C[1]
 =\begin{cases}
 C[1]\otimes B_{k-2}C_{\mathrm{red}}[1]\otimes C[1],&k\geq2,\\
 C[1],&k=1,\\
 \Lambda_{0,nov},&k=0,
 \end{cases}
\tag{4.4.62}\label{eq:4-4-62}
\end{equation}
以及
$\widehat B^{\mathrm{red}}C[1]
=\widehat{\bigoplus}_kB_k^{\mathrm{red}}C[1]$。
\end{definition}

$\widehat d\colon BC[1]\to BC[1]$ 在
$\widehat B^{\mathrm{red}}C[1]$ 上诱导微分
$\widehat d^{\mathrm{red}}$；证明与引理 4.4.24 相同。
定义
$s_{\mathrm{left}},s_{\mathrm{right}}\colon
B_k^{\mathrm{red}}C[1]\to B_{k+1}^{\mathrm{red}}C[1]$：
\[
 s_{\mathrm{left}}(\mathbf x)=e\otimes\mathbf x,
 \qquad
 s_{\mathrm{right}}(\mathbf x)
 =(-1)^{\deg\mathbf x}\mathbf x\otimes e
\]
（$k\geq1$），并令它们在 $B_0^{\mathrm{red}}C[1]$ 上为零。

\begin{lemma}{4.4.63}\label{lem:4-4-63}
在 $\widehat{\bigoplus}_{k\geq1}B_k^{\mathrm{red}}C[1]$ 上，
\[
 \widehat d^{\mathrm{red}}\circ s_{\mathrm{left}}
 +s_{\mathrm{left}}\circ\widehat d^{\mathrm{red}}
 =\widehat d^{\mathrm{red}}\circ s_{\mathrm{right}}
 +s_{\mathrm{right}}\circ\widehat d^{\mathrm{red}}
 =\operatorname{id}.
\]
\end{lemma}

证明只需直接计算，此处从略。选取分裂
\begin{equation}
 C=C_{\mathrm{red}}\oplus\Lambda_{0,nov}[e].
\tag{4.4.64}\label{eq:4-4-64}
\end{equation}
令
$\mathfrak f_{k,\mathrm{red}}\colon
B_k^{\mathrm{red}}C_1[1]\to C_2[1]$（$k=0,1,2,\ldots$）
为一列同态。类似于 \eqref{eq:4-4-28}\translatornote{原书在右端每一因子上都写
$\mathfrak f_{k,\mathrm{red}}(\mathbf x_i)$，而下页说明
$\mathbf x_i\in B_{k_i}C_{1,\mathrm{red}}[1]$；各因子的输入长度
并不必等于块数 $k$，故译文改为各自的 $k_i$。}，对
$\mathbf x_i\in B_{k_i}C_{1,\mathrm{red}}[1]
\subset B_{k_i}C_1[1]$，定义
$\widehat{\mathfrak f}_{\mathrm{red}}\colon
BC_1[1]\to BC_2[1]$：
\[
 \begin{aligned}
 &\widehat{\mathfrak f}_{\mathrm{red}}
 (\mathbf x_1\otimes e\otimes\cdots\otimes e\otimes\mathbf x_k)\\
 &\qquad=
 \mathfrak f_{k_1,\mathrm{red}}(\mathbf x_1)
 \otimes e\otimes\cdots\otimes e\otimes
 \mathfrak f_{k_k,\mathrm{red}}(\mathbf x_k).
 \end{aligned}
\]
\begin{lemma}{4.4.65}\label{lem:4-4-65}
满足下列条件 (4.4.66) 的一列算子
$\mathfrak f_{k,\mathrm{red}}\colon
B_k^{\mathrm{red}}C_1[1]\to C_2[1]$
（$k=0,1,\ldots$）与含幺滤过 $A_\infty$ 同态
$C_1\to C_2$ 一一对应：
\begin{align}
 \widehat{\mathfrak f}_{\mathrm{red}}
 \circ\widehat d^{\mathrm{red}}
 &=\widehat d\circ\widehat{\mathfrak f}_{\mathrm{red}},
 \tag{4.4.66.1}\label{eq:4-4-66-1}\\
 \mathfrak f_{k,\mathrm{red}}\circ s_{\mathrm{left}}
 &=\mathfrak f_{k,\mathrm{red}}\circ s_{\mathrm{right}}=0,
 \tag{4.4.66.2}\label{eq:4-4-66-2}\\
 \mathfrak f_{1,\mathrm{red}}(e)&=e,
 \tag{4.4.66.3}\label{eq:4-4-66-3}\\
 \mathfrak f_{0,\mathrm{red}}&\equiv0
 \pmod{\Lambda_{0,nov}^{+}}.
 \tag{4.4.66.4}\label{eq:4-4-66-4}
\end{align}
\end{lemma}

证明是直接计算，与引理 4.4.30 的证明类似。

现在证明命题 4.4.52 的含幺版本。设已构造
\[
 \widetilde{\mathfrak F}_{k,i}\colon
 B_k^{\mathrm{red}}\overline{\mathfrak C}_1[1]
 \longrightarrow\overline{\mathfrak C}_2[1],
 \qquad k=0,1,2,\ldots,\quad i=0,1,\ldots,n,
\]
它们诱导
\[
 \widetilde{\mathfrak F}^{(n)}_k
 =\sum_{i=0}^{n}\widetilde{\mathfrak F}_{k,i}T^{\lambda_i}
 \colon B_k^{\mathrm{red}}\mathfrak C_1[1]
 \longrightarrow\mathfrak C_2[1],
\]
并满足 \eqref{eq:4-4-66-2}--\eqref{eq:4-4-66-4}。
将其延拓为余代数同态
\[
 \widehat{\widetilde{\mathfrak F}}^{(n),\mathrm{red}}
 \colon\widehat B^{\mathrm{red}}\mathfrak C_1[1]
 \longrightarrow\widehat B\mathfrak C_2[1],
\]
由此诱导含幺同态
\[
 \widehat{\widetilde{\mathfrak F}}^{(n)}
 \colon\widehat B\mathfrak C_1[1]
 \longrightarrow\widehat B\mathfrak C_2[1].
\]
假设它满足 \eqref{eq:4-4-53-1}--\eqref{eq:4-4-53-3}，
目标是构造 $\widetilde{\mathfrak F}_{k,n+1}$。
注意，$n=0$ 的情形已在第 4.4.4 小节证明。

按 \eqref{eq:4-4-54-1}--\eqref{eq:4-4-54-3} 定义
\[
 \begin{gathered}
 \operatorname{Err1}_n\in
 \operatorname{Der}(B\overline{\mathfrak C}_1[1],
 B\overline{\mathfrak C}_2[1]),\\
 \operatorname{Err2}_n\in
 \operatorname{Der}(B\overline C_1[1],
 B\overline{\mathfrak C}_2[1]),\\
 \operatorname{Err3}_n\in
 \operatorname{Der}(B\overline{\mathfrak C}_1[1],
 B\overline C_2[1]).
 \end{gathered}
\]

\begin{lemma}{4.4.67}\label{lem:4-4-67}
对 $*=1,2,3$，算子 $\operatorname{Err*}_n$ 由
$B^{\mathrm{red}}\overline C_1\oplus B_0\overline C_1$
或
$B^{\mathrm{red}}\overline{\mathfrak C}_1
\oplus B_0\overline{\mathfrak C}_1$
上的算子诱导。此外，
$\operatorname{Err*}_n\circ s_{\mathrm{left}}
=\operatorname{Err*}_n\circ s_{\mathrm{right}}=0$。
\end{lemma}

\begin{proof}
对 $\operatorname{Err2}_n$ 和 $\operatorname{Err3}_n$，
结论直接来自定义；对 $\operatorname{Err1}_n$，证明与引理
4.4.35 类似。
\translatornote{原书末句再次写
$\operatorname{Err3}_n$；由于前句已把它列入直接来自定义的情形，
而需仿照引理 4.4.35 处理的是同态的交换误差，故此处改为
$\operatorname{Err1}_n$。}
\end{proof}

利用引理 4.4.67，其余证明与命题 4.4.52 相同。由此得到定理
4.2.34 的含幺滤过版本，定理 4.2.34 全部得证。
